\chapter{Applications}
\label{ch:apps}

I apply the procedure to three datasets chosen to span the range of separability. It
separates little among professional macroeconomic forecasters, finds a top tier among
forecasters in a geopolitical judgement tournament, and separates competent methods from
systematic failures among those submitted to a forecasting competition. The same
\texttt{rankci} pipeline runs on all three; only the loss panel changes. In each case
the width of the confidence sets measures how far the dataset can rank.

\section{Data and common structure}
\label{sec:apps-setup}

Each application supplies a $(\text{observation}\times\text{unit})$ loss panel. Rows are
the observation axis (survey quarters, tournament questions, or held-out months), columns
are the units ranked (forecasters or methods), and each cell is a loss (squared error,
Brier score, or sMAPE). In every case a shared per-row shock is the common shock that
cancels in the pairwise differences: a crisis quarter, a question's realized outcome, an
unusually warm or cold month. The panels differ in balance. The U.S.\ survey panels are
heavily unbalanced, the euro-area and tournament panels are almost balanced, and the
competition panel is balanced.

\section{Macroeconomic forecasters: the null regime}
\label{sec:apps-macro}

Consensus forecasters are genuinely close, and the procedure reports as much: the rank
sets are wide, most of them the trivial full interval, and no best forecaster can be
named. This is the \emph{null} regime, and the ``cannot separate'' verdict is the point.
The data are the U.S.\ Survey of Professional Forecasters (SPF) and its euro-area
counterpart run by the ECB.

\subsection{Why these panels call for dependence-robust inference}
\label{sec:apps-spf-dependence}

Both forms of dependence that motivate the construction (\S\ref{sec:intro-dependence})
are present in these data, and cross-sectional dependence is the dominant one. I measure
them on exactly the panels ranked below: the one-quarter-ahead forecasts of U.S.\ nominal
output, real output and the unemployment rate, and the annual-target forecasts of
euro-area HICP inflation and real GDP growth. The forecast error $e_{t,j}$ is forecast
minus realization, and the loss is $X_{t,j}=e_{t,j}^2$. Appendix~\ref{app:spf-dependence}
gives the definitions, the tests and the full results.

\paragraph{Serial dependence.} The U.S.\ errors are serially correlated. The median lag-1
autocorrelation of an individual forecaster's error is $0.33$--$0.40$, and between $43\%$
and $63\%$ of forecasters have a lag-1 autocorrelation outside the white-noise band,
against about $5\%$ that would be flagged with no serial correlation. The losses inherit
this persistence (median lag-1 autocorrelation $0.23$--$0.39$). The dependence is
short-lived: beyond the first lag the autocorrelations lie inside the band. Overlapping
information sets imply this pattern. A one-quarter-ahead forecast is issued before the
outcome of the current quarter is published, so consecutive errors share an unobserved
shock even if forecasters are fully efficient. For the euro area, at most 27 annual
observations per forecaster reveal no serial correlation (2--3\% of forecasters flagged),
although the test has little power at that length. A standard error that assumed
independence would therefore understate uncertainty on the U.S.\ panels and risk
separating forecasters falsely. The HAC layer of Chapter~\ref{ch:serial} costs little
where the persistence is absent. On the pairwise loss differences the Andrews plug-in
selects a median of only one to three lags, and a single lag on the euro-area panels.

\paragraph{Cross-sectional dependence.} Cross-sectional dependence is much stronger, and
it is present in both surveys. The mean correlation between two forecasters' errors is
$0.76$--$0.97$. Pesaran's CD statistic, approximately standard normal under
cross-sectional independence, lies between $147$ and $301$. A common period effect
explains $82$--$98\%$ of the pooled error variance, against $2$--$3\%$ without a common
component. The dependence is not an artefact of a few crises: rank correlations
($0.67$--$0.91$) and dropping 2020--21 ($0.73$--$0.92$) leave it essentially unchanged.
The losses are correlated almost as strongly ($0.65$--$0.97$). This dependence has two
consequences for the method.

First, the estimated mean losses $\thetahat_j$ are strongly dependent across forecasters.
Treating forecasters as independent cannot approximate the joint distribution of the
contrasts that the simulated critical value of Chapter~\ref{ch:downstream} requires. A
full covariance is needed (Chapter~\ref{ch:cross}).

Second, differencing removes most, but not all, of the common shock. Between $82\%$ and
essentially $100\%$ of the loss variance cancels in $\dser{jk}_t$, as \eqref{eq:cancel}
anticipates. For squared errors, however, the common component sits in the error, not
the loss. Writing $e_{t,j}=f_t+u_{t,j}$ gives
$\dser{jk}_t=(u_{t,j}^2-u_{t,k}^2)+2f_t(u_{t,j}-u_{t,k})$, so the common error $f_t$
re-enters as a scale factor. In the data, the size of the loss differences tracks the
squared consensus error (Spearman correlation $0.63$--$0.89$), and $94$--$100\%$ of the
sum of squared loss differences falls in the tenth of periods with the largest common
error. The difference series are therefore strongly heteroskedastic, with their variance
concentrated in a few episodes.

\paragraph{Implication.} The U.S.\ panels combine this concentrated, non-stationary
variance with heavy imbalance: the eight best-covered forecasters report in only $52\%$
of the quarters. Different pairs are therefore compared on different windows, which
contain different crisis episodes. This is the regime of \S\ref{sec:mc-mechanism}: the
pairwise standard errors are mutually inconsistent, and a level covariance reconstructed
from them loses calibration, whereas the direct construction does not.
Table~\ref{tab:ngdp} below shows exactly this calibration loss on NGDP. The euro-area
panels are almost balanced ($99\%$ coverage) but short, and the same strong common shock
dominates them. The HAC and covariance layers of
Chapters~\ref{ch:serial}--\ref{ch:cross} are built to handle these two forms of
dependence.

Part of the dependence in the U.S.\ output panels is mechanical. These variables are
forecast in levels, so benchmark revisions and base-year changes shift every forecaster's
error at once. Re-expressing the same forecasts as growth rates removes these episodes.
In growth rates the cross-sectional dependence remains strong (mean correlation
$0.74$--$0.77$), while the serial dependence is weaker but still well above chance (29\%
of forecasters flagged); see Appendix~\ref{app:spf-dep-rebasing}.

\subsection{The U.S.\ Survey of Professional Forecasters}
\label{sec:apps-spf}

Across indicators the SPF is close to a null panel, and no indicator names a best
forecaster. The SPF, run quarterly since 1968, collects point forecasts of the main U.S.\
macroeconomic aggregates from professional forecasters. I score each forecaster on each
target quarter by one-quarter-ahead squared error against the Philadelphia Fed real-time
data, keep the eight best-covered forecasters, and rank them with the direct
construction. Table~\ref{tab:spf-indicators} reports the result for three indicators.

The output panels are close to null. On real output only two sets drop a single end
rank. On nominal output three sets are non-trivial, forecaster~65's by two ranks
($[3,8]$) and two others by a single end rank. Unemployment is the most informative
panel: the maximum pairwise $t$-statistic reaches $4.0$, and seven of the eight sets are
non-trivial. Forecaster~65, whose mean loss is the second smallest, is confirmed worse
than five of its seven rivals on their shared quarters and confined to the bottom three
($[6,8]$). Four of the other sets exclude only the last rank, however, and six
forecasters keep rank one in their set.

The MDS route returns the trivial $[1,8]$ for every forecaster on all three panels, the
calibration failure that the NGDP panel below makes visible. The block bootstrap
benchmark (\S\ref{par:mbb}) returns the same sets on real output and unemployment; on
nominal output it is wider by one rank for two forecasters ($[2,8]$ for forecaster~65,
and the trivial $[1,8]$ for forecaster~433). Table~\ref{tab:app-spf-sets} in
Appendix~\ref{app:additional} lists every set.

\tabSpfIndicators

The five best-covered NGDP forecasters show the calibration failure of
\S\ref{sec:mc-mechanism} on real data. Table~\ref{tab:ngdp} ranks them under both
simulation routes. The direct construction separates three of the five, whereas the MDS
route collapses every set to $[1,5]$. On this unbalanced panel the MDS route targets a
different object (Remark~\ref{rem:unbalanced-targets}), and on this heavy-tailed one it
also loses calibration: its studentized draws reach a standard deviation of $12$ and a
two-sided $95\%$ critical value of $24$, against $1.0$ and $2.7$ for the direct
construction.

\tabNgdp

The clearest statement the SPF supports concerns forecaster~65. The $\tau$-best
confidence set, projected from the joint rank sets of the direct construction
(\S\ref{sec:taubest}), makes it precise. Forecaster~65 is excluded from both the top-1
and the top-2 confidence sets, which contain $\{84, 426, 428, 433\}$; only at $\tau=3$
does the set widen to all five. Forecaster~65 has the smallest \emph{marginal} mean
squared error, because its own covered quarters were unusually calm, yet every rival
beats it on their \emph{shared} quarters. The overlap-based contrasts the procedure uses
detect this reversal, where a ranking on marginal means would be fooled.

\subsection{The ECB SPF for the euro area}
\label{sec:apps-ecb}

The ECB SPF for the euro area tells the same story: at these horizons and sample sizes
professional macroeconomic forecasters can hardly be separated, and no best forecaster
can be named (Table~\ref{tab:app-ecb-sets}). I rank the eight best-covered forecasters on
both headline indicators, HICP inflation and real GDP growth, using current-year
forecasts from the first-quarter round. I use the direct construction at $\alpha=0.2$,
as for the U.S.\ panels. On real GDP growth every rank set is the trivial $[1,8]$, and
on HICP inflation only two sets drop a single end rank; the maximum pairwise
$t$-statistic reaches only $2.39$ (HICP) and $1.89$ (RGDP). The panels are almost
balanced, and the MDS route and the block bootstrap give the same sets on both
indicators.

\section{A geopolitical tournament: partial separation}
\label{sec:apps-gjp}

Among tournament forecasters the ranking is \emph{partially} informative: the best
forecaster is confidently in the top three of twenty and the worst is pinned to last,
though the middle remains wide. The Good Judgment Project (GJP) ran a multi-year
tournament in which volunteer forecasters assigned probabilities to geopolitical
questions. I score each forecaster on each question by the canonical daily-averaged
multi-category Brier score. A question's realized outcome is shared by everyone who
answered it, so it is exactly the common shock of \S\ref{sec:apps-setup}; the GJP
literature neutralizes it by standardizing Brier scores within question. I pool the
first two tournament years and restrict to independent individuals and the twenty most
active (212 questions).

The $\tau$-best confidence set for the single best narrows the field to three of twenty.
At the same level, $\alpha=0.2$, the block bootstrap gives the same picture: two of the
twenty sets narrow by one rank, and the normalised width is $0.68$ under the bootstrap
against $0.69$ for the direct construction. The ranking reflects genuine skill, not luck:
the mean Brier scores on two disjoint halves of the questions correlate at a Spearman
$\rho$ of $0.79$ (Figure~\ref{fig:gjp-caterpillar}; Table~\ref{tab:app-gjp} lists every
set and the split-half ranks).

\figGjpCaterpillar

\section{Forecasting methods on the M3 competition: separating failure from competence}
\label{sec:apps-mcomp}

Applied to forecasting \emph{methods} rather than people, the procedure separates
competent methods from systematic failures but names no winner. I do not re-implement
the methods; I rank the forecasts the competing teams actually submitted to the M3
competition, scored against the held-out values the organisers published.

\subsection{The data}
\label{sec:apps-mcomp-data}

The M3 competition \citep{makridakis2000m3} is the third of the large forecasting
competitions organised by Makridakis and co-authors. It contains 3{,}003 real time
series: 645 yearly, 756 quarterly, 1{,}428 monthly and 174 of other frequencies, drawn
from six domains (micro, industry, macro, finance, demographic and other). For each
series the organisers withheld the final observations (6 for yearly, 8 for quarterly and
18 for monthly series), and every participating method had to forecast them from the
remaining history. The organisers then released both the held-out values and the
forecasts of all 24 submitted methods. I use the copy distributed with the R package
\texttt{Mcomp} \citep{hyndman2018mcomp}, which also records a short description of each
series.

Ranking the methods over all 3{,}003 series would pool very different targets: a firm's
monthly shipments and a country's GDP deflator have little in common. I therefore rank
the methods on a single indicator, chosen because it should by nature be forecastable:
the monthly mean air temperature in Paris. M3 contains this record as seventeen monthly
series (N2784--N2800). On inspection they are consecutive, non-overlapping eight-year
blocks of one record that runs from 1857 to 1992. In each block the first 78 months
(January of the first year to June of the seventh) are the estimation window, and the
last 18 months (July of the seventh year to December of the eighth) are held out. The
panel thus has $n=17\times 18=306$ rows, one per (block, horizon) and ordered in calendar
time, and $p=24$ columns, one per method. Every method forecasts every block, so the
panel is balanced.

The loss is the symmetric absolute percentage error used to score the competition,
\begin{equation}
  X_{t,j} \;=\; 200\,\frac{|F_{t,j}-A_t|}{|F_{t,j}|+|A_t|},
  \label{eq:smape}
\end{equation}
where $A_t$ is the realised value and $F_{t,j}$ is method $j$'s forecast. The mean of
$X_{t,j}$ over $t$ is method $j$'s sMAPE. sMAPE is not invariant to a shift of the
measurement origin, and temperature in degrees Celsius is close to zero in winter. The
competition stores the record in a rescaled, strictly positive unit (between $200$ and
$11{,}400$), and I score on that scale exactly as the competition did.

\subsection{The forecasting methods}
\label{sec:apps-mcomp-methods}

\citet{makridakis2000m3} group the 24 submissions into six classes.
Table~\ref{tab:mcomp} gives the class of each method.

\paragraph{Naive and simple methods.} Naive2, the competition's reference benchmark,
removes the seasonal pattern with a classical multiplicative decomposition, carries the
last seasonally adjusted value forward, and reapplies the seasonal index $s$:
$\hat y_{T+h}=(y_T/s_T)\,s_{T+h}$. Single exponential smoothing (Single) updates a level
$\ell_t=\alpha y_t+(1-\alpha)\ell_{t-1}$ and forecasts $\hat y_{T+h}=\ell_T$ at every
horizon.

\paragraph{Explicit trend models.} These are mostly members of the exponential
smoothing family \citep{hyndman2008ets}. Holt's method adds a trend to the level,
\[
  \ell_t=\alpha y_t+(1-\alpha)(\ell_{t-1}+b_{t-1}),\qquad
  b_t=\beta(\ell_t-\ell_{t-1})+(1-\beta)b_{t-1},\qquad
  \hat y_{T+h}=\ell_T+h\,b_T .
\]
The damped-trend method of \citet{gardner1985damped} (Dampen) shrinks the extrapolated
trend by a factor $\phi\in(0,1)$, $\hat y_{T+h}=\ell_T+(\phi+\phi^2+\dots+\phi^h)\,b_T$,
so the forecasts level off at long horizons. Holt--Winters (Winter) adds a
multiplicative seasonal component that is updated together with the level and trend.
Comb S-H-D is the equal-weight average of Single, Holt and Dampen. As for Naive2, the
organisers applied these non-seasonal benchmarks to seasonally adjusted data and
reseasonalised the forecasts. PP-Autocast is a commercial implementation of
damped-trend smoothing, Robust-Trend extrapolates a robustly estimated trend, and
Theta-sm is a variant of the Theta method.

\paragraph{Decomposition.} The Theta method \citep{assimakopoulos2000theta}, the overall
winner of M3, splits the seasonally adjusted series into two ``theta lines'': the
fitted linear trend ($\theta=0$) and a line with twice the local curvature of the data
($\theta=2$). It extrapolates the first linearly and the second by simple exponential
smoothing, and averages the two. \citet{hyndmanbillah2003} show that the result is
simple exponential smoothing with a drift equal to half the slope of the linear trend.

\paragraph{ARIMA-type models.} B-J automatic, the three AutoBox variants, and AAM1 and
AAM2 identify and estimate seasonal ARIMA models
$\phi(L)\Phi(L^{12})(1-L)^d(1-L^{12})^D y_t=\vartheta(L)\Theta(L^{12})\varepsilon_t$
automatically. They differ in their identification rules and, for AutoBox, in how they
treat outliers and interventions. ARARMA first applies a long-memory autoregressive
filter and then fits an ARMA model to the filtered series.

\paragraph{Expert systems.} ForecastPro, ForecastX and SmartFcs are commercial packages
that choose among exponential smoothing and ARIMA models by rules or in-sample criteria.
Rule-based forecasting (RBF; \citealp{collopyarmstrong1992}) combines several
extrapolation methods with weights set by rules that encode features of the series. The
two Flores--Pearce variants are also rule-based expert systems.

\paragraph{Neural network.} Automat ANN fits an automatically specified feed-forward
neural network.

The classes differ in how they handle the annual cycle that dominates a temperature
record. Some model it explicitly (Winter, the seasonal ARIMA models). Some remove it
before extrapolating and restore it afterwards (Naive2, Single, Holt, Dampen, Theta).
The results below show that a few methods do not always capture it at all.

\subsection{Dependence in the panel}
\label{sec:apps-mcomp-dependence}

Both forms of dependence are present. Cross-sectional dependence is weaker than in the
SPF panels; serial dependence in the loss differences is comparable.
Across methods, a month that turns out unusually warm or cold penalises every method
that forecast a typical month; it is the common shock of \eqref{eq:cancel}.
The mean correlation between two methods' losses is $0.65$. For the median pair, the
variance of the difference is $65\%$ smaller than the sum of the two marginal variances,
$1-\Var(\dser{jk}_t)/[\Var(X_{t,j})+\Var(X_{t,k})]=0.65$. Over time, the 18 horizons in
a block all come from one forecast origin, so the losses at adjacent horizons share the
same misjudged level. The median lag-1 autocorrelation of a loss difference is $0.13$,
and the Andrews plug-in selects a median of three lags (range $0$--$19$). Because the
panel is balanced, \S\ref{sec:equivalence} predicts that the \MDS{} and direct
$\OmegaDirect$ constructions coincide up to the effect of the per-pair bandwidths
(Proposition~\ref{prop:equivalence}).

\subsection{Results}
\label{sec:apps-mcomp-results}

The ranking is informative at both ends. Table~\ref{tab:mcomp} and
Figure~\ref{fig:mcomp-caterpillar} report the rank confidence sets at $\alpha=0.1$.
Mean sMAPE ranges from $10.1$ (B-J automatic) to $37.8$ (Flores--Pearce1).

At the top, twelve methods cannot be separated. Their sets all start at rank one, and
they make up the $\tau$-best confidence set for $\tau=1$: B-J automatic, ForecastPro,
ForecastX, Dampen, Comb S-H-D, Theta, Winter, Holt, Single, RBF, SmartFcs and
PP-Autocast. The first eleven of these have mean sMAPEs within one percentage point of
one another ($10.10$--$11.02$), and on this record the M3 winner, Theta, is statistically
indistinguishable from simple exponential smoothing. The other twelve methods are
confidently excluded from first place. For several of them the lower bound is far from
one: Theta-sm cannot rank better than ninth, and ARARMA no better than thirteenth.

At the bottom, three methods are pinned to the last three ranks: Flores--Pearce2
($[22,23]$), Robust-Trend ($[22,24]$) and Flores--Pearce1 ($[23,24]$). The block-level
losses show why (Table~\ref{tab:app-m3-blocks}). Both Flores--Pearce variants issue a
constant forecast, with no seasonal pattern at all, in twelve of the seventeen blocks.
Robust-Trend follows the cycle in some blocks but in eight blocks has an sMAPE above
$25$, more than twice that of the leading methods.

PP-Autocast shows what the procedure does with an unreliable method. Its mean sMAPE
($12.8$) ranks it nineteenth, yet its set is $[1,21]$, the widest in the table. It is
close to the leaders in most blocks, but in one block (1953--60) its sMAPE is $37$,
against about $8$ for the leaders. That single episode makes its loss differences highly
variable, so the procedure can neither exclude it from first place nor place it in the
bottom group. A ranking of the point estimates would place it
nineteenth.

The conclusions do not depend on the covariance construction. The \MDS{} and
$\OmegaDirect$ sets are identical for all but four methods, which differ by one rank, as
the balanced-panel equivalence predicts up to the per-pair bandwidths. The normalised
width is $0.59$ under $\OmegaDirect$ and $0.60$ under \MDS{}. The block bootstrap widens
four sets by one rank, leaves the membership of the top group and the bottom three
unchanged, and has width $0.60$.

Tightening $\alpha$ to $0.05$ leaves the picture essentially unchanged: the same twelve methods share rank one,
and the same three remain the only methods excluded from the top twenty, although
Robust-Trend's set widens to $[21,24]$ and PP-Autocast's to $[1,22]$. Treating each
block as a single observation, with its loss averaged over the 18 horizons ($n=17$),
widens the top group to fourteen methods and raises the best upper bound from $12$ to
$14$. The two Flores--Pearce variants stay pinned at the bottom, and Robust-Trend moves
to $[21,24]$ (Table~\ref{tab:app-m3-robust}).

On a single, inherently forecastable indicator, therefore, the procedure does not pick
out a winner among competent methods, because the data cannot. It does confidently
separate those methods from methods that fail systematically.

\tabMcomp

\figMcompCaterpillar

\section{One procedure, three regimes}
\label{sec:apps-threeregimes}

Under the identical pipeline and at a common level $\alpha=0.1$, the normalised width of
the confidence sets falls from $0.96$ on the macro panel, where only two of the eight
forecasters are separated, each by a single end rank, to $0.73$ on the tournament and
$0.59$ on the competition (Table~\ref{tab:threeregimes}). The two informative panels
are informative in different places. The tournament identifies a top tier and pins one
forecaster exactly. The competition separates a large competent group from a small group
of systematic failures, but pins no method exactly. The procedure does not manufacture
separation that the data cannot support, and it does not hide separation that is
present. The width of the confidence set is a faithful read of how much the data can
rank.

\tabThreeRegimes
